\documentclass[11pt]{article}
\usepackage{amsmath,amsthm,amssymb}
\usepackage[margin=1in]{geometry}
\newtheorem{theorem}{Theorem}
\title{Infrastructure test: an elementary summation identity}
\author{Outreach workflow test fixture}
\date{31 October 2026}
\begin{document}
\maketitle
\begin{abstract}
This synthetic manuscript tests source delivery, PDF compilation and private
review. The identity below is classical; no novelty or research contribution is
claimed. This document is not intended for public publication.
\end{abstract}
\begin{theorem}\label{thm:test}
For every integer $n\geq 1$,
\[
  \sum_{k=1}^{n} k = \frac{n(n+1)}{2}.
\]
\end{theorem}
\begin{proof}
For $n=1$, both sides equal $1$. If the identity holds for an integer $n\geq1$,
then
\[
\sum_{k=1}^{n+1} k
=\frac{n(n+1)}2+(n+1)
=\frac{(n+1)(n+2)}2.
\]
Induction proves the stated identity.
\end{proof}
\paragraph{Test provenance.}
This source was generated by Codex for an explicitly authorized infrastructure
test. It contains no private correspondence and is restricted to private review.
\end{document}
